 Q_1^1(1)=1\ \WIN,\quad Q_2^1(1)=3\ \WIN,\quad
 Q_1^0(1)=2\ \LOSS,\quad Q_2^0(1)=4\ \LOSS.
\]
Hence a boundary \DRAW\ must have accumulated at least one factor of $3$.

If the reached boundary has tail exponent one, put $n=k+1$.  If it has
tail exponent two, its children are the common raw child and
$S_{k+1,1}^e$.  If the raw child is \DRAW, again put $n=k+1$; otherwise the
raw child is \WIN\ and forces $S_{k+1,1}^e$ to be \DRAW, in which case put
$n=k+2$.  Thus in all cases one obtains a boundary \DRAW\ configuration whose
raw/signed exits are governed, for some $n\ge1$, by
\[
  3^n+1-2e=2^jJ(y).
\]
The exact $2$-adic valuation is
\[
\begin{array}{c|cc}
 &n\text{ odd}&n\text{ even}\\ \hline
 e=0&j=2&j=1\\
 e=1&j=1&j=2+v_2(n).
\end{array}
\]
For $j\ge2$ the raw and signed exits form the adjacent factor-free frame
\[
  Q_{j-1}^{1-e}(J(y)),\qquad Q_j^{1-e}(J(y)),
\]
containing a continuing \DRAW; by Definition~\ref{def:entry-controls} this is a boundary entry, with its finite tail length recorded in $\tau$. For $j=1$ one has
\[
  y=\frac{3^{n-1}-1}{2}.
\]
When $e=0$ the raw exit is $A(y)$; when $e=1$ it is $B(y)$. The raw and
signed states are the two exits of the same boundary \DRAW\ configuration.
If the signed exit is \DRAW, it is already the exponent-one form of an
obligation. If the raw exit is the continuing \DRAW\ and the signed exit is
finite, that finite state is necessarily \WIN; choosing a canonical \LOSS\
child of it supplies an exact finite routing witness, hence a loss-sibling
entry. If both exits are \DRAW, use the signed obligation. Thus every row
reaches a typed entry after a finite pre-entry computation; no impossible
``source below zero'' comparison is used.
\end{proof}

\begin{lemma}[$A$-selecting side diamond]\label{lem:a-side}
Let $x>0$, $e=\alpha(x)$, $w=Q_1^e(J(x))$, and $y=A(x)$. Then
\[
  B(w)\in\{A(y),B(y)\}.
\]
\end{lemma}

\begin{proof}
The source-boundary identity gives
$A(w)=Q_1^{1-e}(J(y))=4A(y)+1+e$. For $e=0$ the appended bits are $01$;
for $e=1$ they are $10$. If the first appended bit repeats the last bit of
$A(y)$, deleting the alternating suffix leaves $A(y)$; otherwise the deletion
continues through the maximal alternating suffix of $A(y)$ and leaves
$B(y)$.
\end{proof}

\begin{lemma}[Universal $B$-transfer diamond]\label{lem:b-diamond}
Let $x>0$, let $e=1-\alpha(x)$, put $a=J(x)$ and $g=1-e$, and factor
\[
  3a+1-2e=2^jJ(y),\qquad j\ge2,
\]
so $y=B(x)$. Define
\[
 P=Q_1^e(a),\quad U=Q_2^e(a),\quad
 D=Q_j^g(J(y)),\quad C=Q_{j-1}^g(J(y)),\quad F=A(U).
\]
Then $D=A(P)$ and $C=B(P)=B(U)$. If $X$ is the common child of $D,C$ and
$Y$ is the other child of $C$, then
\[
  B(F)=Y.
\]
Consequently an obligation in the $B$-selecting phase transfers a continuing
\DRAW\ to the adjacent frame $\{D,C\}$; when $j=2$ the transferred frame is
again an obligation in the opposite phase.
\end{lemma}

\begin{proof}
The formulas for $D$ and $C$ are the signed boundary formula and the
shared-child identity. The remaining equality is a direct substitution in the
three cases $j=2$, $j=3$, and $j\ge4$: the constant suffix of $D$ determines
its common child $X$, while the one-step longer factor state $F$ has the same
raw side $Y$. For the outcome assertion, a common child of a \DRAW\ cannot be
\LOSS. If the exponent-one member is \DRAW, the frame $\{D,C\}$ already
contains a \DRAW. In the parent form of an obligation, the only remaining
orientation has the exponent-two member \DRAW; if $C$ were \WIN, outcome
recursion through the displayed diamond would force $Y$ both \LOSS\ and a
child of the \DRAW\ state $F$, impossible. Hence $C$ is \DRAW. For $j=2$,
$D,C$ are exactly the exponent-two/one pair over $J(y)$.
\end{proof}


\section{Audit-hardened routing reduction}\label{sec:routing}

The preceding sections are local arithmetic and outcome lemmas.  The remaining
issue is not another binary identity; it is the global bookkeeping needed to
compose the local diamonds without silently installing an unrelated proof
height or silently changing a ranked source anchor.

\subsection{The old anchor shortcut is permanently withdrawn}

The earlier routing proof attached a retained integer $p$ to an $A$-streak and,
at a later $B$-selecting cursor $y$, used an estimate of the form
\[
 y\le \frac{3p+1}{2}
\]
to conclude that a high-valuation $B$ step lowered the retained source.  That
provenance is false after more than one canonical $A$ continuation.

\begin{lemma}[Anchor stress test]\label{lem:anchor-stress}
Let $x_0=20$ and $x_{i+1}=A(x_i)$.  Then
\[
 x_0,x_1,x_2,x_3=20,30,45,68.
\]
The canonical phases remain $A$-selecting through the first three sources and
switch at $68$.  The selected $B$ transition at $68$ has valuation $3$ and
\[
 B(68)=25>20.
\]
Consequently no proof may infer a decrease relative to the original anchor
$20$ merely from the later high-valuation $B$ transfer.
\end{lemma}

\begin{proof}
Directly,
\[
 A(20)=30,\qquad A(30)=45,\qquad A(45)=68,
\]
and
\[
 \alpha(20),\alpha(30),\alpha(45),\alpha(68)=1,0,1,1.
\]
The lifted phase flips at every canonical $A$ continuation, so the required
phase at $68$ is $0$, opposite to $\alpha(68)=1$.  Moreover
\[
 3J(68)+1=616=2^3\cdot77,
 \qquad 77=J(25),
\]
which gives the asserted valuation and $B(68)=25$.
\end{proof}

The revised routing architecture therefore treats the displayed integers in an
$A$-streak as \emph{temporary cursors}.  A cursor is never promoted to a ranked
source merely because it is smaller than a later cursor.  A control reset may
be justified only by one of the two certificates below:
\begin{enumerate}[label=(R\arabic*)]
\item an actual DRAW state whose retained coefficient source is proved smaller
than the currently ranked source by an exact factor-free/source lemma;
\item a finite WIN/LOSS proof token that is proved to be a proper proof-tree
descendant of a token already present in the ranked multiset (or the unique
initial token installation).
\end{enumerate}

\subsection{Finite control is not semantic completeness}

The finite control vocabulary consists of the following names:
\begin{center}
\begin{tabular}{lll}
\texttt{boundary\_entry} & \texttt{loss\_sibling\_entry} & \texttt{a\_obligation}\\
\texttt{a\_test\_4} & \texttt{a\_test\_3} & \texttt{a\_test\_2}\\
\texttt{a\_test\_1} & \texttt{b\_select} & \texttt{b2\_first}\\
\texttt{b2\_ready} & \texttt{factor\_fork} & \texttt{high\_return}\\
\texttt{marked\_tail} & \texttt{short\_lift} & \texttt{terminal\_macro}.
\end{tabular}
\end{center}
The finite graph of these names is only a bookkeeping device.  A graph edge is
admissible in the mathematical proof only when its guard is derived from the
actual game states and its finite witnesses have the exact parent/child
incidence claimed for them.  The accompanying JSON file therefore calls its
edges \emph{declared routing edges}, not proved game transitions.

The arithmetic part of a $B$-selecting transfer is safe without any retained
anchor comparison.  Lemma~\ref{lem:b-diamond} gives an adjacent frame over the
actual selected source $z=B(y)$.  If the selected valuation is $2$, this frame
is again an exact obligation.  If the valuation is at least $3$, it is routed
to the factor-frame normalizer.  In neither row do we assert that $z$ is below
an older pre-streak source.

\section{The two remaining proof obligations}\label{sec:remaining}

The following statements are deliberately separated from the proved lemmas.
They are the only assertions needed to turn the present reduction back into an
unconditional no-DRAW proof, and they are \emph{not} claimed proved in this
version.

\paragraph{Obligation P1: seeded $A/B_2$ reset provenance.}
Suppose the inner routing layer already contains a ranked finite proof token.
Follow an actual DRAW through any finite canonical $A$-streak and then through
a valuation-two $B$ transfer that returns to an obligation.  Before the
control is reset to a new \texttt{a\_obligation}, one must exhibit either
(i) a genuine ranked source decrease, or (ii) a finite witness which is a
strict proof-tree descendant of an already ranked token.  The statement must
include multi-$A$ streaks and all four exceptional residue classes.  It is not
enough that the new source is WIN, and it is not enough that it is smaller
than the immediate cursor.

\paragraph{Obligation P2: semantic exhaustiveness of the factor router.}
For every actual DRAW entering any factor-router control, the complete
outcome split must be derived from the game relation.  The relevant controls
are \texttt{factor\_fork}, \texttt{high\_return},
\texttt{marked\_tail}, \texttt{short\_lift}, and
\texttt{terminal\_macro}.  Every case must land either in a declared
non-reset routing edge or in a strict source/token edge.  In particular, every reset to a higher control
must carry the explicit descendant certificate for the token it replaces.
No finite endpoint of unrelated height may be inserted merely because it is
WIN or LOSS.

The long proof ledger in the repository contains many of the local identities
needed for P1 and P2, including the factor-fork and high-return descendants.
They are useful evidence, but the present manuscript does not promote the two
global composition statements to theorems until their guards and provenance
are checked as one closed transition system.

\begin{proposition}[Conditional routing closure]\label{prop:conditional-routing}
Assume P1 and P2.  Then the marked routing relation admits a well-founded
lexicographic/multiset ranking in which every control reset is preceded by a
strict retained source or certified proof-token decrease, while all remaining
routing is finite by the bounded control graph and its explicit tail/factor
counters.
\end{proposition}

\begin{proof}
The proof is the standard well-founded-fibre assembly.  Rank finite proof
witnesses by the natural multiset sum
\[
 \Phi(M)=\mathop{\#}_{h\in M}\omega^h
\]
from Lemma~\ref{lem:token-rank}.  A certified descendant replacement strictly
lowers $\Phi$.  A genuine retained-source edge lowers the numerical source and
may reset later data.  At fixed retained source and token multiset, P1 and P2
say exactly that no edge which resets the finite control is unranked.  Delete
the strict source/token edges.  What remains is the declared finite control
graph together with same-control recurrences whose explicit natural tail or
factor counter decreases.  The residual relation is well-founded.  The
well-founded-fibre lemma then gives the full marked relation.
\end{proof}

\section{Conditional no-DRAW consequence and present status}

\begin{theorem}[Conditional no-DRAW theorem]\label{thm:conditional}
If Obligations P1 and P2 hold, then the target no-DRAW statement of
Section~\ref{target:nodraw} follows.
\end{theorem}

\begin{proof}
Assume a DRAW exists and choose a least coefficient source among DRAW
positions.  The long-tail reduction, the four explicit exponent-one/two
boundary cases of Proposition~\ref{prop:boundary-reduction}, and the
factor-free base entry of Proposition~\ref{prop:base-entry} yield either a
strict smaller-source DRAW (contradicting minimality) or a typed entry.  The
factorful exponent-one row is handled by Proposition~\ref{prop:factorful-entry};
no reverse-factor lemma is used at exponent one and no coefficient/source
shortcut is used with a factor $3^k$.

Under P1 and P2, Proposition~\ref{prop:conditional-routing} makes every
outcome-compatible marked continuation well-founded.  But every DRAW has no
LOSS child and at least one DRAW child, so the local outcome diamonds could be
iterated forever while retaining an actual DRAW member.  This contradicts
well-foundedness.  Hence no DRAW exists in the conjugated game, and
Proposition~\ref{prop:first-normal} transfers finite remoteness to every
original positive odd start.
\end{proof}

\paragraph{Status.}
The implication above is proved; P1 and P2 are not.  Accordingly version 5.0
\textbf{does not claim that the prize problem is solved}.  This wording is
intentional.  It prevents a finite control certificate, a regression test, or
a moving repository supplement from being mistaken for the missing semantic
proof.

\section{Verification boundary and reproducibility}

The package contains three kinds of checks.
\begin{itemize}[leftmargin=*]
\item \texttt{verify\_v5\_0.py} checks the exact arithmetic identities used in
the proved part, the old source/coefficient counterexample, the permanent
$20\to30\to45\to68$ anchor counterexample, and the declared finite control
graph.
\item \texttt{routing\_certificate\_v5\_0.json} records controls and guards.
Every edge is labelled either \texttt{pure-routing}, \texttt{strict-source},
\texttt{strict-token}, or \texttt{OPEN-P1/P2}.  The verifier rejects an
unlabelled reset edge.
\item \texttt{OPEN\_PROOF\_OBLIGATIONS\_v5\_0.md} gives the exact mathematical
targets still required for an unconditional theorem.
\end{itemize}

These checks are deliberately narrower than a formal proof.  No end-to-end
Lean formalization is claimed.

\appendix
\section{Permanent adversarial regression cases}\label{app:regressions}

The following two examples are normative negative tests for every future
revision.

\paragraph{Coefficient/source separation.}
For $s=1$, $J(1)=5$, $a=3J(1)=15$, and $e=1$, the common boundary child is
$R(44)=22$.  Its canonical coefficient is $11<15$, but $11=J(3)$, so the
coefficient source grows from $1$ to $3$.  Therefore a smaller canonical
coefficient does not imply a smaller source when powers of three are present.

\paragraph{Pre-streak anchor separation.}
The calculation of Lemma~\ref{lem:anchor-stress} gives
$20\to30\to45\to68$ and $B(68)=25>20$.  Therefore a high-valuation $B$
transfer after several $A$ steps does not by itself lower the source retained
at the beginning of the streak.

\section{Declared routing inventory}\label{app:routing}

The table is a checklist, not a proof of P1/P2.
\begin{center}
\small
\begin{tabular}{@{}p{.25\linewidth}p{.61\linewidth}@{}}
\toprule
Control & Allowed non-strict successor classes \\
\midrule
\texttt{boundary\_entry} & finite tail countdown; then
  \texttt{a\_obligation} or \texttt{factor\_fork} \\
\texttt{loss\_sibling\_entry} & finite tail countdown; then
  \texttt{factor\_fork} \\
\texttt{a\_obligation} & bounded $A$ test or factor fork; a smaller-source
  DRAW may re-enter only through a proved source comparison \\
\texttt{a\_test\_4,3,2,1} & next lower test or \texttt{b\_select} \\
\texttt{b\_select} & valuation $2$ to \texttt{b2\_first}; valuation
  $\ge3$ to \texttt{factor\_fork}; no retained-anchor descent is asserted \\
\texttt{b2\_first} & \texttt{b2\_ready} \\
\texttt{b2\_ready} & reset requires P1 \\
\texttt{factor\_fork} & \texttt{high\_return} or strict certified exit \\
\texttt{high\_return} & reset requires the P2 descendant certificate \\
\texttt{marked\_tail} & shorter tail/short lift, or strict P2 token exit \\
\texttt{short\_lift} & \texttt{a\_obligation} \\
\texttt{terminal\_macro} & every reset requires a P2 token certificate \\
\bottomrule
\end{tabular}
\end{center}

\begin{thebibliography}{9}
\bibitem{Guy1986}
R. K. Guy, John Isbell's Game of Beanstalk and John Conway's Game of
Beans-Don't-Talk, \emph{Mathematics Magazine} 59(5) (1986), 259--269.

\bibitem{Guy1996}
R. K. Guy, Unsolved Problems in Combinatorial Games, Problem 42, in
\emph{Games of No Chance}, MSRI Publications 29, Cambridge University Press,
1996.

\bibitem{Althofer2024}
I. Alth\"ofer, M. Hartisch, T. Zipproth, Analysis of a Collatz Game and Other
Variants of the $3n+1$ Problem, \emph{Advances in Computer Games} (2024),
123--132.

\bibitem{AlthoferPage}
I. Alth\"ofer, Collatz Problems: The Two-Player $3n+-1$ Game,
\url{https://althofer.de/collatz-prizes.html}.

\bibitem{Repo}
G. Pochuev, Optimal $3n\pm1$ Game: Proof and Verification Repository,
\url{https://github.com/Grisha-Pochuev/3n-plus-minus-1-game}.
\end{thebibliography}

\end{document}
